\documentclass[10pt,a4paper]{amsart}
\usepackage[margin=19mm]{geometry}
\usepackage{amsmath,amssymb,mathtools}
\usepackage[scr=rsfs,cal=euler]{mathalfa}
\usepackage{xfrac}
\usepackage[unicode,hidelinks,bookmarks=false]{hyperref}
\newcommand{\ie}{i.e.\ }
\newcommand{\cf}{cf.\ }
\newcommand{\resp}{respectively\ }
\newcommand{\Subsection}{\subsection}
\providecommand{\nobreakdash}{\nobreak\hbox{-}}
\setlength{\emergencystretch}{2em}
\multlinegap0pt
\usepackage{bm}
\usepackage{bbm}
\usepackage{dsfont}
\usepackage{xspace}
\usepackage{cancel}
\usepackage{mathtools}
\usepackage{cleveref}
\usepackage[shortlabels]{enumitem}
\usepackage{amsthm}
\makeatletter
\@ifundefined{defi}{\newtheorem{defi}{Defi}}{}
\@ifundefined{lem}{\newtheorem{lem}{Lem}}{}
\@ifundefined{thm}{\newtheorem{thm}{Thm}}{}
\makeatother
\newcommand{\At}{\operatorname{At}}
\newcommand{\A}{\mathfrak{A}}
\newcommand{\B}{\mathfrak{B}}
\newcommand{\C}{\mathfrak{C}}
\newcommand{\D}{\mathfrak{D}}
\newcommand{\EPFL}{\exists^{+}\!\mathsf{FL}}
\newcommand{\FL}{\mathsf{FL}}
\newcommand{\at}{\operatorname{at}}
\newcommand{\false}{\bot}
\newcommand{\qr}{\operatorname{qr}}
\newcommand{\tp}{\operatorname{tp}}
\newcommand{\true}{\top}
\title[Preservation and definability for the fluted fragment]{Preservation and definability for the fluted fragment}
\author{Yiwen Ding}
\date{Source: arXiv:2607.12970v2 (CC BY 4.0)}
\begin{document}
\maketitle

\noindent\emph{Source and licence.} Preservation and definability for the fluted fragment. Version arXiv:2607.12970v2 (CC BY 4.0), distributed under the Creative Commons Attribution 4.0 International licence, as recorded in the arXiv licence metadata for this exact version and shown on the abstract page https://arxiv.org/abs/2607.12970v2. Full rights notice: licenses/arxiv-2607.12970-CC-BY-4.0.txt.\par\smallskip

\noindent\textit{Editorial scope.} Counted unit: the Thm whose source label is \texttt{thm:finite-hpt} in the article (source0.tex lines 1178--1188, source label \texttt{thm:finite-hpt}), delivered together with the article's own complete original proof of it (source0.tex lines 1190--1261, one \texttt{proof} environment of 72 lines). No mathematics is written, repaired or re-derived: statement and proof are the article's own text line for line. Required context retained uncounted: lem:type-finite (lines 523--526); def:characteristic-formula (lines 551--601); eq:characteristic-formula (lines 562--570); lem:chi-basic (lines 607--629); lem:canonical-query (lines 749--758); lem:reset (lines 842--848); lem:ep-hom (lines 1001--1012). Every definition, display and label of those lines is retained unchanged. Omitted from this package and not redistributed: 1--522, 527--550, 571--606, 630--748, 759--841, 849--1000, 1013--1177, 1189--1189, 1262--2089. Nothing inside a retained excerpt is omitted or altered: every retained body is delivered exactly as the article's lines are written, so no omission receipt of any kind accompanies this package. Citations: the retained excerpts carry no citation command at all, so no bibliography entry of the article is transcribed and no reference is redistributed. Reference adaptation: none was needed and none was made; every reference inside the delivered unit resolves inside it, so no unresolved reference or citation is produced. Printed slips kept literal: no printed word, symbol, hypothesis or number of the counted statement or of its proof was altered. Layout and class: the article's own class is \texttt{lmcs} with packages including amssymb, array, booktabs, hyperref, cleveref; the wrapper is the lane's pinned \texttt{amsart} editorial wrapper at 10pt on A4, which restates the theorem-like environments the retained text needs and adds the article's own macro definitions that the retained text needs (\texttt{\textbackslash A}, \texttt{\textbackslash At}, \texttt{\textbackslash B}, \texttt{\textbackslash C}, \texttt{\textbackslash D}, \texttt{\textbackslash EPFL}, \texttt{\textbackslash FL}, \texttt{\textbackslash at}, \texttt{\textbackslash false}, \texttt{\textbackslash qr}, \texttt{\textbackslash tp}, \texttt{\textbackslash true}), with extra wrapper packages (bm, bbm, dsfont, xspace, cancel, mathtools, cleveref, enumitem). The delivered numbering is the unit's own. Assets: no retained span contains a figure environment, a graphics-inclusion command, a PDF-page inclusion, a table or a TikZ picture; no image file is copied into this package, the package declares no asset dependency and no image file is redistributed with it. Licence: version arXiv:2607.12970v2 of this article is distributed under the Creative Commons Attribution 4.0 International licence, as recorded in the arXiv licence metadata for this exact version and shown on the abstract page https://arxiv.org/abs/2607.12970v2. A full rights notice is delivered at licenses/arxiv-2607.12970-CC-BY-4.0.txt. Characters: every retained excerpt is pure ASCII, so the delivered engine, XeLaTeX with its default Latin Modern text font, reports no missing character.
\begin{lem}[Finiteness of type codes]\label{lem:type-finite}
For every $m,k\geq0$, there are only finitely many rank-$k$ type codes at
level $m$.
\end{lem}

\begin{defi}[Characteristic formula]\label{def:characteristic-formula}
For every $m,k\geq0$ and every rank-$k$ type code $t$ at level $m$, define
its \emph{characteristic formula}
$\chi_t^{m,k}\in\EPFL[m](\sigma)$ by induction on $k$, simultaneously for all
levels $m$.  At rank $0$, we have $t\subseteq\At_m$, and we put
\[
  \chi_t^{m,0}:=\bigwedge_{\alpha\in t}\alpha,
\]
where, throughout this definition, empty conjunctions are interpreted as
$\true$.  For the induction step, let $t=(p,S)$ be a rank-$(k+1)$ type code
at level $m$, and put
\begin{equation}
  \chi_t^{m,k+1}
  :=
  \left(\bigwedge_{\alpha\in p}\alpha\right)
  \wedge
  \bigwedge_{s\in S}
     \exists x_{m+1}\chi_s^{m+1,k}.
  \label{eq:characteristic-formula}
\end{equation}
By \cref{lem:type-finite},~\eqref{eq:characteristic-formula} is a finite
conjunction, so the
recursion is well-defined.

In~\eqref{eq:characteristic-formula}, each conjunct of the form
$\exists x_{m+1}\chi_s^{m+1,k}$ has a separate quantifier.  Different
successor type codes may therefore use different witnesses, even though the
same variable symbol $x_{m+1}$ is reused.  All these conjuncts are evaluated
under the same assignment to $x_1,\ldots,x_m$, although only a suffix of these
variables may actually occur freely.

For a particular structure and tuple, abbreviate
\[
  \chi_{\A,\bar a}^{k}
  :=
  \chi_{\tp_{\A}^{k}(\bar a)}^{m,k},
  \qquad |\bar a|=m.
\]
For every $\sigma$-structure $\A$ and every $n\geq0$, define the
\emph{rank-$n$ positive characteristic sentence} associated with the empty
tuple by
\[
   \Theta_{\A}^{n}
   :=
   \chi_{\A,\epsilon}^{n}
   \in\EPFL[0](\sigma).
\]
Notice that, by construction,
$\chi_{\A,\bar a}^{k}\in\EPFL[m](\sigma)$ and
$\qr(\chi_{\A,\bar a}^{k})\leq k$ whenever $|\bar a|=m$.
\end{defi}

\begin{lem}\label{lem:chi-basic}
For every $m,k\geq0$, every $\sigma$-structure $\A$, and every tuple
$\bar a\in A^m$:
\begin{enumerate}[label=\textup{(\roman*)}]
  \item $\A\models\chi_{\A,\bar a}^{k}[\bar a]$;
  \item for every $\sigma$-structure $\B$ and every $\bar b\in B^m$, if
  $\tp_{\A}^{k+1}(\bar a)=\tp_{\B}^{k+1}(\bar b)$, then the following
  conditions hold:
  \begin{itemize}
    \item \emph{Atomic harmony.}
    $\at_{\A}(\bar a)=\at_{\B}(\bar b)$.
    \item \emph{Forth condition.}  For every $a'\in A$, there exists $b'\in B$ such
    that
    $\tp_{\A}^{k}(\bar a a')=\tp_{\B}^{k}(\bar b b')$.
    \item \emph{Back condition.}  For every $b'\in B$, there exists $a'\in A$ such
    that
    $\tp_{\A}^{k}(\bar a a')=\tp_{\B}^{k}(\bar b b')$.
  \end{itemize}
\end{enumerate}
In particular, for every $\sigma$-structure $\A$ and every $n\geq0$,
$\Theta_{\A}^{n}\in\EPFL[0](\sigma)$,
$\qr(\Theta_{\A}^{n})\leq n$, and $\A\models\Theta_{\A}^{n}$.
\end{lem}

\begin{lem}[Canonical-query property]\label{lem:canonical-query}
For every $\sigma$-structure $\B$,
\[
   \B\models\Theta_{\A}^{n}
   \quad\Longleftrightarrow\quad
   \text{there is a homomorphism }\C_{\A}^{n}\longrightarrow\B.
\]
In particular, $\C_{\A}^{n}\models\Theta_{\A}^{n}$.

\end{lem}

\begin{lem}[Disjoint-copy reset lemma]\label{lem:reset}
There is a rank-$n$ fluted bisimulation $(Z_0,\ldots,Z_n)$ between
$\A^{\dagger}$ and $\D$ with $(\epsilon,\epsilon)\in Z_n$.  Consequently,
\[
   \A^{\dagger}\equiv_n^{\FL}\D.
\]
\end{lem}

\begin{lem}\label{lem:ep-hom}
Let $h:\A\to\B$ be a homomorphism of $\sigma$-structures.  For every
$\eta\in\EPFL[m](\sigma)$ and every $\bar a\in A^m$, write
$h(\bar a)=(h(a_1),\ldots,h(a_m))$.  Then
\[
  \A\models\eta[\bar a]
  \quad\Longrightarrow\quad
  \B\models\eta[h(\bar a)].
\]
In particular, every sentence in $\EPFL[0](\sigma)$ is preserved under
homomorphisms.
\end{lem}

\begin{thm}[Finite homomorphism preservation]\label{thm:finite-hpt}
Let $\varphi$ be a sentence in $\FL[0](\sigma)$, and put
$n:=\qr(\varphi)$.  The following are
equivalent, with preservation and equivalence taken over finite
$\sigma$-structures:
\begin{enumerate}[label=\textup{(\roman*)}]
  \item $\varphi$ is preserved under homomorphisms;
  \item $\varphi$ is equivalent to some $\psi\in\EPFL[0](\sigma)$ with
  $\qr(\psi)\leq n$.
\end{enumerate}
\end{thm}

\begin{proof}
For \textup{(ii)}$\Rightarrow$\textup{(i)}, let $\psi$ witness
\textup{(ii)} and let $h:\A\to\B$ be a homomorphism between finite structures.
If $\A\models\varphi$, then $\A\models\psi$, hence $\B\models\psi$, and
therefore $\B\models\varphi$ by \cref{lem:ep-hom}.

For \textup{(i)}$\Rightarrow$\textup{(ii)}, suppose that $\varphi$ is
preserved under homomorphisms between finite $\sigma$-structures.  If $n=0$,
then $\varphi$ is equivalent to $\true$ or $\false$.  If $\varphi$ is
unsatisfiable over finite structures, it is equivalent to $\false$ there.  We
may therefore assume that $n\geq1$ and that $\varphi$ is satisfiable over
finite structures.

As in the proof over all structures, it remains to establish characteristic
entailment and then collect the finitely many root type codes.  We also need
to check that every structure in the transfer argument remains finite.

Fix a finite $\A\models\varphi$, put $\C:=\C_{\A}^{n}$, let $\D$ be the
disjoint union of $n+1$ copies of $\C$, and put
$\A^{\dagger}:=\A\uplus\D$.  Every structure in the transfer chain
\[
   \A\longrightarrow\A^{\dagger}
   \equiv_n^{\FL}\D
   \longrightarrow\C
\]
is finite: $\C$ is finite by construction, $\D$ has finitely many copies, and
$\A^{\dagger}=\A\uplus\D$.  The first and last arrows are the canonical
injection and fold homomorphisms, while the middle equivalence is
\cref{lem:reset}.  Since $\qr(\varphi)=n$, the chain gives
$\C\models\varphi$ using preservation only between finite structures.

If a finite $\B$ satisfies $\Theta_{\A}^{n}$, then
\cref{lem:canonical-query} gives a homomorphism $\C\to\B$.  Both structures
are finite, so homomorphism preservation yields $\B\models\varphi$.  Thus
$\Theta_{\A}^{n}$ entails $\varphi$ over finite structures.

Now let
\[
  \mathcal T_{\varphi}^{\mathrm{fin}}
  :=
  \bigl\{
    \tp_{\A}^{n}(\epsilon)
    \mid
    \A\text{ is a finite $\sigma$-structure satisfying }\varphi
  \bigr\}.
\]
By \cref{lem:type-finite} and finite satisfiability of $\varphi$,
$\mathcal T_{\varphi}^{\mathrm{fin}}$ is finite and non-empty.  Define
\[
  \psi_{\varphi}^{\mathrm{fin}}
  :=
  \bigvee_{t\in\mathcal T_{\varphi}^{\mathrm{fin}}}\chi_t^{0,n}.
\]
By \cref{def:characteristic-formula}, each disjunct belongs to
$\EPFL[0](\sigma)$ and has quantifier rank at most $n$.  Hence
$\psi_{\varphi}^{\mathrm{fin}}\in\EPFL[0](\sigma)$ and
$\qr(\psi_{\varphi}^{\mathrm{fin}})\leq n$.

If a finite $\B$ satisfies $\varphi$, then
$\tp_{\B}^{n}(\epsilon)\in\mathcal T_{\varphi}^{\mathrm{fin}}$, and
\cref{lem:chi-basic}(i) gives
$\B\models\chi_{\tp_{\B}^{n}(\epsilon)}^{0,n}$.  Hence
$\B\models\psi_{\varphi}^{\mathrm{fin}}$.

Conversely, if a finite $\B$ satisfies $\psi_{\varphi}^{\mathrm{fin}}$, then
$\B\models\chi_t^{0,n}$ for some
$t\in\mathcal T_{\varphi}^{\mathrm{fin}}$.  Choose a finite
$\A\models\varphi$ with $t=\tp_{\A}^{n}(\epsilon)$.  Since
$\chi_t^{0,n}=\Theta_{\A}^{n}$, the characteristic entailment proved above
gives $\B\models\varphi$.  Thus $\varphi$ and
$\psi_{\varphi}^{\mathrm{fin}}$ are equivalent over finite structures.
\end{proof}

\end{document}
